\subsection{Monday Lecture}

\begin{prop}[Continuity of a Measure]
   Any measure \( \mu (\Omega) < \infty  \) satisfies the following properties: 
   \begin{enumerate}
       \item[(i)] (Finite additivity): For a finite sequence of disjoint sets \( {E}_{1}, {E}_{2}, \dots, {E}_{k} \)
            \[  \mu \Big(  \sum_{ i=1  }^{ k  } {E}_{i} \Big) = \sum_{ i=1  }^{ k  } \mu({E}_{i}). \]
   \end{enumerate}
\end{prop}
\begin{proof}
Suppose we have an \( \mathcal{F}-   \)algebra where \( \mu : \mathcal{F} \to [0,\infty) \cup  \{ \infty  \}  \). If \( {E}_{n} \) where \( n \in \N  \) is a sequence of disjoint sets; that is, \( {E}_{i} \cap {E}_{j} = \emptyset  \) if \( i \neq j  \). Then 
\[  \mu \Big(  \bigcup_{ i=1  }^{ \infty  }  {E}_{i} \Big) = \sum_{ i=1  }^{ \infty  } \mu({E}_{i}). \]
Note that in (i), given the above, we can write
\[  \mu \Big(  \bigcup_{ i= 1  }^{ k }  {E}_{i} \Big) = \mu \Big(  \sum_{ i=1  }^{ k } {E}_{i} \Big) = \sum_{ i=1  }^{ k } \mu({E}_{i}).  \]
This statement can be easily proven using property (1) of the definition of a measure where all indices \( i \) past \( k \), we can set \( {E}_{i} = \emptyset \). Indeed, we will consider three cases:   
\begin{enumerate}
    \item[(1)] \( i \neq j  \) where \( i,j \in \{  1,2,\dots, k  \}  \) implies \( {E}_{i} \cap {E}_{j} = \emptyset \) where \( {E}_{1}, {E}_{2}, \dots, {E}_{k} \) are all disjoint.
    \item[(2)] \( i \neq j  \) where \( i,j  \in \{  k + 1 , k + 2, \dots  \}  \) implies \( {E}_{i} \cap {E}_{j} = \emptyset \cap \emptyset = \emptyset \).
    \item[(3)] \( i \in \{  1,2,\dots, k  \}  \) and \( j \in \{  k + 1 , k + 2 , \dots  \}  \) implies \( {E}_{i} \cap {E}_{j} = {E}_{i} \cap \emptyset  = \emptyset\).
\end{enumerate}
From (i)-(iii), we have \( {E}_{i}  \) where \( i \in \N  \) is a disjoint sequence of sets. The above proposition includes \( {E}_{n}  \) for \( n \in \N  \) disjoint and 
\[  \mu\Big(\bigcup_{ i = 1  }^{ \infty  }  {E}_{n}\Big) = \sum_{ i=1  }^{ \infty  } \mu({E}_{n}), \]
which implied that \( \bigcup_{ i=1  }^{ \infty  }  {E}_{i} \in \mathcal{F} \). Note that 
because 
\begin{align*}
    \mu \Big(  \bigcup_{ i=1  }^{ \infty  }  {E}_{i}  \Big) &= mu \Big(  \bigcup_{ i =1  }^{ \infty  }  {E}_{i} \Big) = \sum_{ i=1  }^{ \infty  } \mu({E}_{i}) \\
                                                            &= \sum_{ i=1  }^{ k  } \mu({E}_{i}) + \sum_{ i = k + 1  }^{ \infty  } \mu(\emptyset) \\
                                                            &= \sum_{ i=1  }^{ k  } \mu({E}_{i}) + 0  \\
                                                            &= \sum_{ i=1  }^{ k  } \mu({E}_{i}).
\end{align*}
\end{proof}

Recall that if \( f  \) is continuous, then \( {x}_{n} \to x  \) which implies that \( f({x}_{n}) \to f(x) \).

\begin{prop}[Measure is Continuous from Below]
    If \( {E}_{n}  \) is a non-decreasing sequence to \( E  \), \( E \in \mathcal{F} \) and \( {E}_{n} \in \mathcal{F} \) for all \( n \in \N \), then \( \mu({E}_{n})  \) converges to \( \mu(E) \) where \( \mu(E) \in \{ \sup E  , \infty   \}  \). 
\end{prop}

\subsection{Wednesday Lecture}

Redo the following propositions on homework.

\begin{prop}[B.2.5 (Redo some of the details in your proof)]
   \( \lim_{ x \to y }  f(x) = v  \) if and only if \( \lim_{ n \to \infty  }  x^{(n)}  = y \) and for all \( n \in \N  \), \( x^{(n)} \neq y  \), then  
   \begin{enumerate}
       \item[(*)] \( \lim_{ n \to \infty  }  x^{(n)}  = y  \) where \( x^{(n)} \neq y  \) for all \( n \in \N  \) means for all \( \epsilon > 0  \), there exists an \( N \in \N  \) such that for all \( n > N  \) (\( n \in \N  \)), \( d(x^{(n)}, y) < \epsilon \).
        \item[(*)] If \( \lim_{ x \to y }  f(x) = v  \) means that for all \( \epsilon > 0  \), there exists \( \delta > 0  \) such that if \( {d}_{A}(x,y) < \delta \), then \( {d}_{B}(f(x), v) < \epsilon \).
   \end{enumerate}
\end{prop}

\begin{proof}
\( (\Longrightarrow) \) Since \( \lim_{ x \to y }  f(x) = v  \). Let \( \epsilon > 0  \) be given. Then there exists \( \delta > 0  \) such that if \( {d}_{A}(x,y) < \delta \), then
\( {d}_{B}(f(x), v) < \epsilon \). Then if \( \lim_{ x \to \infty  }  x^{(n)}  = y  \) (\( x^{(n)} \neq y   \)) for all \( n \in \N  \).

\begin{center}
 Our goal is to show that for all \( \epsilon > 0  \), there exists an \( N \in \N  \) such that for all \( n > N  \), \( {d}_{B}(f(x^{(n)}), v) < \epsilon \).   
\end{center}
  Then for the aforemention \( \delta > 0  \), there exists \( N \in \N  \) if \( n > N  \), \( d(x^{(n)},y) < \delta \), then \( {d}_{B}(f(x^{(n)}), v) < \epsilon \). Since \( \epsilon  \) was arbitrary, for all \( \epsilon > 0  \), there exists an \( N \in \N  \) such that if \( n > N  \), then \( {d}_{B}(f(x^{(n)}), v) < \epsilon \). Hence, \( \lim_{ n \to \infty  }  f(x^{(n)}) = v  \).

  (\( \Longleftarrow \)) Suppose \( \lim_{ n \to \infty  }  x^{(n)} = y  \) where (\( x^{(n)} \neq y  \) for all \( n \in \N \)) and for sake of contradiction \( \lim_{ x \to y }  f(x) \neq v  \); that is, there exists an \( \epsilon > 0  \) such that for all \( \delta > 0  \), there exists an \( x \neq y  \in A  \) such that if \( {d}_{A}(x,y) < \delta \), \( {d}_{B}(f(x), v) \geq \epsilon \). For our fixed \( \epsilon > 0  \), let \( \delta = 1   \). Then there exists an \( {x}_{1} \neq y \in A  \) such that if \( {d}_{A}(x^{(1)}, y) < 1 \), then \( {d}_{B}(f({x}_{1}), v) \geq \epsilon \). Then for \( \delta = \frac{ 1 }{ 2 }   \), there exists an \( {x}_{2} \neq y  \in A  \) such that if \( {d}_{A}(x^{(2)}, y) < \frac{ 1 }{ 2 } \), then \( {d}_{B}(f(x^{(3)}), v) \geq \epsilon \). We continue in this manner until, we have a sequence \( ({x}_{n}) \subseteq  A  \) where \( {x}_{n} \neq y  \) for all \( n \in \N  \) where if \( {d}_{A}(x^{(n)}, v) < \frac{ 1 }{ n }  \), then \( {d}_{B}(f(x^{(n)}), v) \geq \epsilon \).  
\end{proof}

\begin{prop}
   Prove that if \( d(x^{(n)},y) < \frac{ 1 }{ n }  \), then \( \lim_{ n \to \infty  }  x^{(n)} = y  \).
\end{prop}

\begin{prop}[Redo]
   Let \( f: A \to B  \) and \( g : B \to C  \) be continuous functions on \( A  \) and \( B  \), respectively. Then \( g \circ f : A \o C   \) is continuous on \( A  \).
\end{prop}

\begin{proof}
Our goal is to show that \( {B}_{n} \) for \( n \in \N  \) is a disjoint sequence of sets. 

Note that \( {B}_{1} = {E}_{1} \) (\textbf{This is a homework problem}) and \( {B}_{n} = {E}_{n} \setminus  {E}_{n-1} \). Let \( i \neq j  \) where \( i,j \in \{  2,3,\dots \}  \). Without loss of generality, let \( i < j  \). Then our goal is to show that \( {B}_{i} \cap {B}_{j} = \emptyset \). Note that \( {B}_{i} = {E}_{i} \setminus  {E}_{i-1} \) and \( {B}_{j} = {E}_{j} \setminus  {E}_{j-1} \).  Then 
\begin{align*}
{B}_{i} \cap {B}_{j} &= ({E}_{i} \cap {E}_{j-1}^{c} ) \cap  ({E}_{j} \cap {E}_{j-1}^{c}) \\ 
                     &= ({E}_{i} \cap {E}_{j}) \cap ({E}_{i-1}^{c} \cap {E}_{j-1}^{c}). \tag{*}
\end{align*}
Since \( {E}_{i} \subseteq  {E}_{j} \), \( {E}_{i} \cap {E}_{j} = {E}_{i} \) and since \( {E}_{i-1} \subseteq {E}_{j-1}    \) implies that \( {E}_{j-1}^{c} \subseteq  {E}_{i-1}^{c}\), and so \( {E}_{i-1}^{c} \cap {E}_{j-1}^{c} = {E}_{j-1}^{c} \). Then we have (*) implies that \( {B}_{i} \cap  {B}_{j} =   {E}_{i} \cap {E}_{j-1}^{c} \). If \( x \in {E}_{i} \cap {E}_{j-1}^{c} \) implies \( x \in {E}_{i} \) and \( x \notin {E}_{j-1} \) but \( {E}_{i} \subseteq  {E}_{j-1} \) (\textbf{Polish the rest of this for homework}). 

\end{proof}

\begin{prop}[(3) Continuity from Above]
    
\end{prop}


